% Generated draft: requires mathematical and source review.
\documentclass[11pt]{article}
\usepackage[margin=25mm]{geometry}
\usepackage{fontspec}
\setmainfont{Latin Modern Roman}
\usepackage{amsmath,amssymb,mathtools,unicode-math}
\setmathfont{Latin Modern Math}
\usepackage{enumitem,needspace,xurl,hyperref}
\setlength{\parindent}{0pt}
\setlength{\parskip}{5pt}
\setlength{\emergencystretch}{3em}
\newcommand{\N}{\mathbb{N}}
\newcommand{\Z}{\mathbb{Z}}
\newcommand{\Q}{\mathbb{Q}}
\newcommand{\R}{\mathbb{R}}
\newcommand{\C}{\mathbb{C}}
\newcommand{\F}{\mathbb{F}}
\newcommand{\A}{\mathbb{A}}
\newcommand{\PP}{\mathbb{P}}
\newcommand{\E}{\mathbb{E}}
\newcommand{\eps}{\varepsilon}
\newcommand{\dd}{\,\mathrm{d}}
\begin{document}
\section*{UnsolvedMath ID: 30001152}
\textbf{Source problem code:} catalogue entry
\subsection{Definitions and mathematical statement}
A membrane model must specify its geometric domain, its mass distribution, the relevant differential operator and boundary conditions, and the admissible variations. For such a model \(X\), write \(m(X)\) for its total mass. When its Neumann spectrum is discrete and has an ordered sequence of positive eigenvalues, write
\[
0<\mu_1(X)\leq\mu_2(X)\leq\cdots,
\]
listing eigenvalues with multiplicity and omitting all zero eigenvalues. Write \(\sigma_k(X)\) for its \(k\)th Steklov eigenvalue once the Steklov problem and its indexing convention have been specified. The source explicitly says ``nonzero'' for the Neumann eigenvalue but does not specify whether the Steklov indexing includes the zero eigenvalue.

The supplied question is: for prescribed membrane mass, how large can \(\mu_k\) and \(\sigma_k\) be? Its conditional optimization formulation is the following. For a prescribed mass \(M>0\) and index \(k\in\{1,2,\ldots\}\), after choosing admissible membrane classes \(\mathcal A_N\) and \(\mathcal A_S\) on which the respective eigenvalues are defined, determine
\[
\sup_{\substack{X\in\mathcal A_N\\m(X)=M}}\mu_k(X)
\qquad\text{and}\qquad
\sup_{\substack{X\in\mathcal A_S\\m(X)=M}}\sigma_k(X).
\]
These classes are placeholders for missing source specifications, not additional hypotheses asserted by this draft. No finiteness or attainment of either supremum is asserted.

The supplied record does not specify the admissible domains or geometries, the mass functional, density assumptions, regularity, connectedness, or the precise weighted or unweighted Neumann and Steklov boundary value problems. In particular, it does not say whether mass is an interior quantity for the Neumann problem and a boundary quantity for the Steklov problem. Consequently, the displayed optimization problems cannot be made uniquely precise from the supplied record alone; their model, admissible classes, and Steklov indexing require clarification from the source.
\subsection{Short English statement}
For a membrane with prescribed mass, determine the largest possible \(k\)th positive Neumann eigenvalue and \(k\)th Steklov eigenvalue. The admissible membrane models, meaning of mass, and Steklov indexing are unspecified, so the optimization formulation remains conditional. The supplied record labels the problem open; this draft does not independently verify that status.
\subsection{Sources}
\begin{itemize}
\item Supplied local UnsolvedMath registry record, permanent ID \(30001152\), exact source problem code \texttt{catalogue entry}, title: \emph{Mass-Constrained Neumann and Steklov Eigenvalue Maximization}. Collection: \emph{Oberwolfach Reports Open Problems}; supplied proposed year: 2009. No individual proposer is identified.
\item Supplied source citation: \emph{Low Eigenvalues of Laplace and Schrödinger Operators} (2009), \emph{Oberwolfach Reports}, pp.\ 355--428, DOI: \url{https://doi.org/10.4171/owr/2009/06}. The record describes Oberwolfach Reports as documenting mathematical workshops and their open-problem discussions; no problem-specific page is supplied.
\item Supplied literature source, labeled ``Steklov shape-optimization survey'': \url{https://doi.org/10.1007/s00039-021-00573-5}. No further bibliographic details are supplied.
\item Literature assessment dated exactly \texttt{2026-08-21}: the supplied curation labels the problem open and states that sharp low-index inequalities are known in several classes, while the general higher-eigenvalue maximization problem under a mass constraint is not settled. This is supplied dated curation derived from the OWR record, not independently verified literature assessment; its time-sensitive status and cited claims require independent verification.
\item The registry labels the statement \texttt{corrected\_verified} and reports that its mathematical content was checked against the official Oberwolfach report, with notation, context, and question form normalized without changing the source-level claim. This is a supplied verification claim, not verification performed for this draft.
\end{itemize}
\end{document}
